\section{Introduction}
\label{sec:intro}

Recent multi-turn tool-agent work \citep{coevolve2026,multiturnguide2025} uses
group-relative policy-gradient (PG) methods such as GRPO
\citep{deepseekmath2024}. These methods can
learn from outcome verification alone, but their token-level credit assignment is
coarse: the same trajectory-level advantage is broadcast across long sequences
of reasoning, API calls, tool observations, and final answers. Self-distillation
(SD) offers an attractive complement: use the policy's own successful behavior,
or a teacher conditioned on privileged training-time information, to supply a
denser signal \citep{sdpo2026,opsd2026,skillsd2026,sdar2026}. The hard part is
making that signal help the policy gradient rather than fight it.

\paragraph{The hidden failure.} We first show, on the corrected $\tau^3$-airline
customer-service task set in the $\tau$-bench family
\citep{taubench2024,tau2bench2025,tau3bench2026}, that naive SD
\emph{can degrade} long-horizon tool-use. Over training, the model
increasingly solves easy information-only cases without using tools, while
state-changing action tasks collapse to zero success (\S\ref{sec:diagnosis}).
The aggregate score conceals the mechanism: the policy is not learning a better
agent; it is abandoning executable tool-use behavior the base or supervised
checkpoint already possessed.

\paragraph{Design implication: use distillation for credit, not as a competing loss.}
The contrast with a $\tau^3$ Vanilla GRPO control using the benchmark
reference-KL is informative: it preserves
tool-use and reaches held-out mean Pass$^1$ $0.583$, whereas SDPO reaches
$0.317$ and collapses on state-changing action tasks. Our subsequent
AppWorld probes led to the
same design constraint. Direct SD pressure can be live and non-inert, yet still
degrade action quality when the teacher target is misaligned with what the
verifier rewards; our diagnostic direct/uncapped SD ablation
falls below the reported Vanilla GRPO control. This evidence motivates
repurposing
teacher/student divergence as a \emph{credit-weighting signal}: it determines
where the broadcast policy-gradient advantage should receive bounded
token-level emphasis, while the
reward-grounded policy optimization remains the update driver. On AppWorld,
this detached signal replaces Vanilla GRPO's reference-KL operationally, but
affects the actor only through credit weights rather than an equivalent penalty.

\paragraph{The method.} We introduce \textbf{Sibling-Guided Credit Distillation
(SGCD)}. Dynamic sampling produces mixed groups of verified successful and failed
sibling rollouts for the same task; \creditreferenceproviderphrase{} summarizes their
contrast into a stepwise credit reference (what succeeded, which branches failed,
where failures deviated). The policy is scored twice on the rollout: once under
the ordinary student context and once under the credit reference as a stop-gradient
teacher. Dense top-$K$ reverse-KL and entropy convert the divergence into
clipped token-level advantage weights inside the policy-gradient surrogate
(\S\ref{sec:method}). No distillation loss touches the actor: distillation guides
credit, and policy gradient performs the update.

\paragraph{Training-time boundary.} The external LLM, sibling evidence, and
credit reference are training-only; the inference-time student sees the same clean
task prompt as the policy, with no oracle, sibling rollouts, external LLM, or credit
reference. Because the divergence signal is detached from the actor parameters
and only reshapes the verified advantage, SGCD removes the teacher-matching
actor-gradient path that can pull the model toward actions the clean deployment
prompt cannot reproduce.

\paragraph{Contributions.}
\begin{enumerate}
  \item A tool-use diagnosis: naive self-distillation can reinforce benchmark
    shortcuts, preserving easy information-only behavior while losing
    state-changing action capability (\S\ref{sec:diagnosis}).
  \item \textbf{SGCD}: a sibling-guided credit-distillation method that builds a
    stepwise credit reference from mixed successful/failed rollouts and uses
    dense teacher/student divergence to reweight, not replace, the PG update
    (\S\ref{sec:method}).
  \item An empirical study on AppWorld and $\tau^3$-airline: AppWorld carries the
    main positive held-out result, while $\tau^3$-airline provides a directional
    held-out comparison and a failure-mode diagnostic. Ablations give lower
    point estimates when the sibling credit reference or divergence-driven
    weighting is removed, and direct SD pressure falls below the reported Vanilla
    GRPO control (\S\ref{sec:experiments}).
\end{enumerate}
